\documentclass[11pt,a4paper]{article}
\usepackage[margin=2.3cm]{geometry}
\usepackage{amsmath}
\usepackage{mathptmx}
\usepackage{booktabs}
\usepackage{tabularx}
\usepackage{xcolor}
\usepackage{caption}
\usepackage{enumitem}
\usepackage[round,authoryear]{natbib}
\usepackage[hidelinks]{hyperref}
\captionsetup{font=small,labelfont=bf}
\setlist{itemsep=2pt,topsep=4pt}
\setlength{\parskip}{5pt}
\setlength{\parindent}{0pt}
\newcolumntype{Y}{>{\raggedright\arraybackslash}X}

\title{\textbf{Wrapping up the soil-moisture foundation-model study}\\[3pt]
\large What was done, where it sits in the literature, the paper's arguments,\\
cheap additions, and the next paper}
\author{Prajwal Khanal}
\date{11 October 2026}

\begin{document}
\maketitle

\section{What was done}

\paragraph{Data.} 990 stations from ISMN \citep{dorigo2021ismn}, ICOS and AmeriFlux with at
least three years of daily data, 2016--2025. Hourly in-situ soil moisture (SM) was QC-filtered
(G/D* flags), averaged to daily values and binned to three depths: 0--10, 10--30 and 30--100\,cm
(885, 682 and 551 stations; 5.14\,M station-days). Splits by station: 573 train, 89 validation,
224 out-of-sample (OOS) stations. Splits by time: 2016--2022 train, 2023--2025 out-of-time
(OOT). OOST is the joint case (unseen station, unseen years). Normalisation statistics were fit
on training stations and years only.

\paragraph{Inputs.} Per station, a 2.24\,km tile:
\begin{itemize}
\item frozen TerraMind \citep{jakubik2025terramind} layer-12 tokens (14$\times$14 at 160\,m) for
  Sentinel-2 L2A (272\,k cloud-screened scenes), Sentinel-1 RTC (497\,k scenes), DEM and LULC;
\item ERA5-Land, 18 variables, 365-day window;
\item SIF (last 50 values), GRACE TWSA (12 months), gridded soil properties;
\item a light 20\,m CNN path on the latest S2/S1 scene (added in \S48).
\end{itemize}

\paragraph{Model.} A single 3-layer transformer (26.7\,M trainable parameters) over 1{,}038
tokens: three depth query tokens, static tokens, satellite history tokens and daily driver tokens.
Huber loss at the station pixel; ERA5 day-dropout 0.15 and modality dropout (ERA5 0.5,
satellite 0.3, fine path 0.2). Selected checkpoint: epoch 10, run
\texttt{baseline\_selected\_20261005} (runbook \S59/\S66).

\begin{table}[h]
\centering\small
\caption{Baseline skill, station-mean ubRMSE (m$^3$\,m$^{-3}$) and Pearson $R$ per depth.
Source: \texttt{eval\_output/baseline\_selected\_20261005/metrics\_summary.csv}.}
\begin{tabular}{l c ccc ccc}
\toprule
 & & \multicolumn{3}{c}{ubRMSE} & \multicolumn{3}{c}{$R$} \\
\cmidrule(lr){3-5}\cmidrule(lr){6-8}
Split & $n$ (0--10) & 0--10 & 10--30 & 30--100 & 0--10 & 10--30 & 30--100 \\
\midrule
Validation (new stations) & 89  & .0529 & .0476 & .0470 & .753 & .748 & .711 \\
OOS (new stations)        & 221 & .0523 & .0483 & .0471 & .722 & .708 & .634 \\
OOT (seen stations, 2023--25) & 362 & .0490 & .0433 & .0437 & .763 & .768 & .720 \\
OOST (new stations, 2023--25) & 114 & .0535 & .0474 & .0495 & .733 & .741 & .688 \\
\bottomrule
\end{tabular}
\end{table}

\begin{table}[h]
\centering\small
\caption{Input reliance on 221 OOS stations: ubRMSE increase when one input group is swapped
with another station's (cross) or with another date at the same station (within). Source:
\texttt{\dots/baseline\_selected\_20261005\_ablation/ablation\_summary.csv}.}
\begin{tabular}{l ccc}
\toprule
Input swapped & 0--10 & 10--30 & 30--100 \\
\midrule
All satellite tokens, cross & +36.8\,\% & +37.1\,\% & +38.4\,\% \\
ERA5-Land, within station   & +39.0\,\% & +30.4\,\% & +13.5\,\% \\
ERA5-Land, cross            & +30.3\,\% & +24.6\,\% & +15.1\,\% \\
Soil properties, cross      & +7.6\,\%  & +17.6\,\% & +25.7\,\% \\
DEM, cross                  & +12.6\,\% & +9.7\,\%  & +18.4\,\% \\
LULC, cross                 & +6.1\,\%  & +6.6\,\%  & +12.7\,\% \\
S1 history, cross           & +4.8\,\%  & +7.7\,\%  & +14.4\,\% \\
S2 history, cross           & +4.9\,\%  & +6.8\,\%  & +8.7\,\% \\
S1 history, within station  & 0.0\,\%   & +1.0\,\%  & 0.0\,\% \\
SIF, cross                  & +1.3\,\%  & +1.6\,\%  & +1.8\,\% \\
20\,m fine path, cross      & +0.5\,\%  & +1.3\,\%  & +0.1\,\% \\
TWSA, cross                 & 0.0\,\%   & 0.0\,\%   & +0.1\,\% \\
\bottomrule
\end{tabular}
\end{table}

\begin{table}[h]
\centering\small
\caption{Negative and null results already measured.}
\begin{tabularx}{\linewidth}{@{}l Y@{}}
\toprule
Test & Result \\
\midrule
LST as auxiliary target (\S53/\S56) & Hurts SM: best val ubRMSE 0.0537 with LST vs 0.0487 without; $R$(0--10) .694 vs .761. \\
20\,m fine path (\S56.2, \S67) & Unused for SM: val 0.0487 intact, 0.0489 zeroed, 0.0496 shuffled. \\
Landsat LST (\S29) & Tile pattern is static (month-to-annual $r=+0.967$); pooled $r=+0.167$ but within-station $-0.077$ (Simpson's paradox). \\
ECOSTRESS DTR (G0) & DTR is the daytime pattern ($r=+0.816$); day vs night $+0.059$. Thermal arm closed. \\
Effective resolution (\S70) & $L_\mathrm{eff}=2240$\,m at 0--10\,cm; within-tile contrast $\sim$4$\times$ below observed. \\
Patchwise model (\S35, \S56) & OOT beat validation (0.0434 vs 0.0507), best epoch 2: site memorisation. Tile-pair sign test 0.50. \\
Depth, wd, regularisation (\S60--\S65) & 1/2/3/6 layers: 0.0510/0.0494/0.0488/0.0500; wd 0.3: 0.0493; static dropout 0.0513; spectral indices 0.0511. None beats the 3-layer baseline. \\
TerraMind layer (\S27b) & kNN skill for station-mean SM: L12 9.8\,\%, L9 6.9\,\%, L6 3.6\,\%, L3 4.0\,\%; tabular covariates 23.4\,\%. \\
\bottomrule
\end{tabularx}
\end{table}

\section{Literature and the gap}

\paragraph{Geospatial foundation models (GFMs).} Prithvi-EO-2.0 \citep{szwarcman2024prithvi},
SatMAE \citep{cong2022satmae}, Galileo \citep{tseng2025galileo}, Copernicus-FM
\citep{wang2025copernicusfm}, DOFA \citep{xiong2024dofa}, AnySat \citep{astruc2024anysat},
Presto \citep{tseng2023presto} and TerraMind \citep{jakubik2025terramind} are pretrained
self-supervised encoders. AlphaEarth \citep{brown2025alphaearth} and TESSERA
\citep{feng2025tessera} release precomputed embeddings. Benchmarks
\citep{lacoste2023geobench,marsocci2024pangaea,wang2025copernicusfm} are dominated by
classification and segmentation; PANGAEA finds GFMs do not consistently beat supervised UNets,
and \citet{xu2025specialized} show tuned supervised baselines match specialised foundation models.
\citet{rolf2024position} argue for EO-specific, spatially aware evaluation, and SatCLIP
\citep{klemmer2025satclip} shows that location alone predicts many regression targets.

\paragraph{GFM embeddings for continuous variables.} The closest SM study
(\citealp{kontogiorgakis2026sm}, preprint; 113 European stations) finds Prithvi embeddings add nothing
over hand-crafted S1/S2/ERA5 features ($R^2$ 0.515 vs 0.514). An IBM abstract fine-tuning Prithvi on
TxSON \citep{chavda2024ibm} reports no fine-tuning numbers. AlphaEarth embeddings are highly stable
between years ($r=0.963$; \citealp{rahman2026aef}, preprint), i.e.\ largely static. A buffered spatial
validation halves the embedding gain \citep{ratsakatika2026buffer}, cross-country crop yield gives no
FM advantage \citep{adjei2026yield}, and frozen encoders underperform for biomass while precomputed
embeddings do better \citep{sialelli2026agb}. A 2026 review \citep{yu2026ecohydro} notes that
pretraining corpora lack microwave and thermal data and that embeddings are mostly static while
ecohydrology needs dynamics. \citet{fang2026alphaearth} (Section~3) is the clearest positive case.

\paragraph{ML soil moisture from in-situ networks.} SoMo.ml \citep{o2021somo} trains an LSTM on
ISMN to 50\,cm at 0.25$^\circ$; GLASS SM \citep{zhang2023glass} and GSSM1km \citep{han2023gssm}
produce 1\,km surface products; \citet{zhang2023glass} is one of few to report random,
site-independent and year-independent validation side by side. Deep-learning comparisons
\citep{fang2017lstm,wang2024hess} use hand-crafted inputs. S1/SMAP downscaling
\citep{das2019smap,bauer2019s1} is surface-only and degrades under dense vegetation.

\paragraph{Validation.} Random CV overstates skill for spatial prediction
\citep{ploton2020spatial,meyer2022maps}; the area of applicability (AOA) of
\citet{meyer2021aoa} has been applied to in-situ SM upscaling \citep{yu2025aoa}.

\paragraph{Gaps this study fills.}
\begin{enumerate}[label=G\arabic*.]
\item No test of a multimodal any-to-any GFM (TerraMind) for soil moisture; existing tests use
  Prithvi only.
\item No GFM soil-moisture study below the surface; we predict 10--30 and 30--100\,cm.
\item Existing GFM soil-moisture studies use 40--113 stations in one region; we use $\sim$990
  stations across climates.
\item No study separates the embedding as a site descriptor from the embedding as a dynamic signal.
\item New-station, new-year and joint generalisation are rarely reported together.
\item Whether the embedding gain survives a station holdout is untested for soil moisture.
\item No fusion of GFM embeddings with reanalysis, SIF and GRACE in one temporal model.
\item GFM benchmarks lack a regression task on a dynamic variable at point sites.
\item Frozen GFM, precomputed embeddings and raw features are not compared on the same SM task.
\item No AOA or embedding-distance uncertainty for GFM-based SM at unseen stations.
\end{enumerate}
Gaps G1--G5 and G7 are covered by work already done. G6, G9 and G10 need the additions in
Section~5. G8 is the next paper.

\section{Fang et al.\ (2026) and what we take from it}

\paragraph{What they do.} AlphaEarth annual 64-d embeddings at 10\,m are masked to corn or soybean
(CDL) and averaged per county-year (6{,}350 records, 992 Corn Belt counties, 2017--2024). An
RBF-SVR is the only model. Validation is leave-one-year-out plus a roll-forward forecast
(train $\le t-1$). Corn: mean $R^2=0.825$, RMSE 14.8\,bu/ac; soybean $R^2=0.814$. Baselines:
training-years county mean ($R^2\approx0.69$) and NDVI+PRISM+GDD (0.687 corn, 0.732 soy).
Diagnostics: Moran's $I$ of relative errors 0.32--0.49 every year (spatially clustered residuals,
read as missing slow drivers); PCA axes explain $<$3\,\% of error variance (variance $\neq$ task
relevance); test-fold permutation importance on the 64 dimensions; a label-subsampling curve.

\begin{table}[h]
\centering\small
\caption{Fang et al.\ (2026) vs this study.}
\begin{tabularx}{\linewidth}{@{}l Y Y@{}}
\toprule
 & Fang et al.\ (2026) & This study \\
\midrule
Target & County yield, once per year & Daily SM at a point, 3 depths \\
Embedding & AlphaEarth, annual, 64-d, precomputed & TerraMind L12 tokens per scene, frozen \\
Other inputs & None & ERA5-Land, SIF, TWSA, soil \\
Model & RBF-SVR & 3-layer transformer \\
Spatial holdout & None (every county in every training fold) & 89 val + 224 OOS stations \\
Temporal holdout & LOYO, roll-forward & OOT 2023--25, OOST \\
Naive baseline & County mean, NDVI+weather & Not yet (Section~5, E1) \\
Residual diagnostics & Moran's $I$, LISA & Not yet (E3) \\
\bottomrule
\end{tabularx}
\end{table}

\paragraph{How it relates.} Both studies ask whether GFM embeddings are analysis-ready features.
Their setting cannot separate site identity from signal: every county is in training, the county
mean alone gives $R^2\approx0.69$, and an annual embedding has no within-season dynamics. Our
result fills exactly that hole: under a station holdout the embeddings carry place, and the
time-varying skill comes from the drivers (Table~2). We cite Fang et al.\ as the positive case and
position our paper as the test it lacks.

\paragraph{What we borrow.} (i) A per-station climatology baseline, the analogue of their county
mean. (ii) Moran's $I$ / LISA on OOS station residuals. (iii) A label-subsampling curve. (iv)
AlphaEarth embeddings as a drop-in static token, which also gives a GFM-vs-GFM comparison (G9).

\section{Paper storyline and arguments}

\paragraph{Working title.} \emph{Frozen foundation-model embeddings describe the site, weather
drives the dynamics: daily root-zone soil moisture at 990 stations.}

\paragraph{One-paragraph story.} Frozen TerraMind embeddings fused with reanalysis and
gravimetric drivers predict daily soil moisture to 1\,m at unseen stations with ubRMSE
0.047--0.052\,m$^3$\,m$^{-3}$, close to stations seen in training. Swap tests show the embeddings act
as a static landscape descriptor that conditions the weather-to-soil-moisture response; they do not
carry the temporal signal. Reliance shifts from weather to soil and terrain with depth. Finer
resolution and thermal inputs do not help. This agrees with recent reports that GFM embeddings are
static, and adds the first station-holdout, multi-depth evidence for soil moisture.

\paragraph{Arguments, each with its number.}
\begin{enumerate}[label=A\arabic*.]
\item \textbf{Transfer to new sites is close to seen sites.} OOS ubRMSE .052/.048/.047 vs
  validation .053/.048/.047 and OOT .049/.043/.044 (Table~1).
\item \textbf{Embeddings carry place, not process.} Swapping all satellite tokens across stations
  costs +37--38\,\%; swapping S1 history within a station costs 0--1\,\%. L12 tokens are nearly
  static across seasons (\S72.2). Consistent with \citet{rahman2026aef,yu2026ecohydro}.
\item \textbf{Dynamics come from the drivers.} ERA5 swap within station +39\,\% at 0--10\,cm,
  falling to +13.5\,\% at 30--100\,cm.
\item \textbf{Reliance inverts with depth.} Soil properties go from +7.6\,\% to +25.7\,\%,
  DEM to +18.4\,\% at 30--100\,cm.
\item \textbf{More resolution and thermal do not help.} Fine path 0.1--1.3\,\%; LST target
  worsens SM by 10\,\%; effective resolution 2.2\,km, the tile size.
\item \textbf{Validation design decides the verdict.} The patchwise model looked good until OOT
  beat validation, a memorisation signature. We report new-station, new-year and joint splits
  together \citep{ploton2020spatial,meyer2022maps}.
\end{enumerate}

\paragraph{Wording rule.} The swap tests measure \emph{reliance}, not \emph{added value}: a model
trained without an input could compensate. Say ``the model relies on'' until the drivers-only run
(E2) exists.

\paragraph{Open bookkeeping.} OOS counts differ: 224 in the split, 222 evaluated, 221 at 0--10\,cm.
The ablation file's reference ubRMSE (.0500/.0479/.0476) uses a different aggregation from Table~1,
so quote only relative changes from it. Fix both before submission.

\section{Cheap additions for this paper}
Ranked by cost. Eval-only items reuse the existing predictions.

\begin{table}[h]
\centering\small
\begin{tabularx}{\linewidth}{@{}l Y l Y@{}}
\toprule
 & Addition & Cost & Why \\
\midrule
E1 & Naive baselines: station climatology (OOT only), persistence, ERA5-only gradient boosting & CPU, eval & Reviewers will ask; we have none. Mirrors Fang's county mean. \\
E2 & Drivers-only model (no satellite tokens), 10 epochs & 1 training run & Turns ``reliance'' into ``added value''. The most important addition. \\
E3 & Moran's $I$ / LISA on OOS residuals & CPU, eval & Shows whether errors cluster in space (missing slow drivers). \\
E4 & AOA: embedding distance to training stations vs OOS error & CPU, eval & Uncertainty for new sites (G10). \\
E5 & Group Shapley (\S73, already planned) & GPU, eval & Clean attribution per input group. \\
E6 & Label-subsampling: train on 25\,\% and 50\,\% of stations & 2 runs & How much data the embeddings save. Optional. \\
E7 & AlphaEarth annual embedding as a static token & GEE pull + 1 run & GFM-vs-GFM (G9). Optional. \\
\bottomrule
\end{tabularx}
\end{table}

Minimum for submission: E1, E2, E3. E4--E7 strengthen it but are not required.

\section{The next paper}
\begin{enumerate}[label=N\arabic*.]
\item \textbf{A soil-moisture benchmark for GFMs.} ISMN stations, fixed new-station / new-year /
  joint splits, three depths. Encoders: TerraMind, Prithvi, AlphaEarth, TESSERA, Galileo, raw
  bands. This fills G8 and reuses the pipeline that already exists.
\item \textbf{Is ``static'' a pretraining limit?} Fine-tune the encoder, or use time-resolved
  embeddings (Presto/Galileo-style pixel time series), and test whether the within-station swap
  effect rises above zero.
\item \textbf{Wall-to-wall maps with uncertainty.} The Overijssel run (\S68) with TWENTE validation
  and an AOA mask.
\item \textbf{Management signals.} Irrigation at Mead (\S69): the irrigation signal appears only
  in dry years, and S2 does not see it.
\end{enumerate}

\bibliographystyle{plainnat}
\bibliography{paper_wrapup}

\end{document}
